\section{Appendix: Scope and Qualifications}
\label{sec:scope}

{\small
\textbf{Edition, formulary and occurrence.} The study treats the Mass
\latin{Dominica vigesima prima post Pentecosten}, II classis, of the \work{Missale
Romanum}, editio typica, Typis Polyglottis Vaticanis 1962, pp.~413--415,
nos.~1699--1708, in the universal 1962 calendar and no particular calendar. Its
occurrence on 18~October 2026, its precedence over the feast of St Luke, the
commemoration of that feast in low and conventual Masses and its omission in a
sung Mass that is not conventual, the oration \latin{pro Fidei propagatione} added
in every Mass on the penultimate Sunday of October, the colour, the Credo and the
Trinity Preface rest on the Missal's own calendar and general rubrics (RG 16, 91,
108, 111, 127~b; RGMR 434, 437~a, 444--446, 453, 475~a, 483, 494~b). Whether a local
Ordinary has appointed another day for the missions, and any diocesan, titular,
dedication or institute calendar or local external solemnity, are local facts not
applied here. The added orations are named by incipit and locus and not
reproduced; they are not part of the formulary.

\textbf{Text control.} Every Latin form of the ten proper elements was read on page
images of the CMAA facsimile of the typical edition and collated with the Pustet
Missal of 1862, in which every element's wording already stands; the collation,
the conclusion forms and the comparison with the Clementine Vulgate are recorded
in \texttt{propers/verified.md}. No witness read outside liturgical books carries
the Introit's antiphon in the Missal's form, the Missal's opening of the Communion
or its added \latin{Domine Deus meus}, or the wording of the Offertory's narrative;
the Gradual's \latin{a generatione et progenie} stands in Cassiodorus's lemma as
transcribed, and the Gospel's \latin{suffocabat} in St Jerome's lemma in Migne. No
chant manuscript, Old Latin or Roman psalter was consulted. The Missal prints the
Offertory antiphon alone; the verses that older chant books sang after it are not
part of this formulary.

\textbf{English.} The scriptural English is the Douay--Rheims in Challoner's revision
at the canonical verses. Where the Missal departs from the Vulgate (the Introit's
antiphon, the Epistle's opening, the Gradual's \latin{a generatione et progenie}, the
Gospel's opening, the Offertory's narrative, and the Communion's shortened and
altered verses) that English does not render the Missal's Latin; the Offertory and
Communion are given in the English of the 1861 Philadelphia hand missal, which also
supplies the opening sentence of the Introit and the three orations. None of these
is an approved liturgical translation of the 1962 Missal. The English of the
Fathers, Doctors and commentators is that of the published translations named in
the References. Where the Latin of St Jerome, St Hilary, St Thomas or Cassiodorus is printed, the
English beside it is a gloss of that Latin and not a translation from a named
edition.

\textbf{Reception.} Each appointed passage was searched in the Latin and Greek
Fathers, the later Doctors and saints, and the liturgical commentators available to
this production, and the witnesses read are those named in the References at their
exact loci. St Hilary's commentary on Ps~118 and St Jerome's on Matthew and
Ephesians were read on the page images of Migne's \work{Patrologia Latina}, and the
exposition of the psalms printed under St Athanasius's name on the page images of
Migne's \work{Patrologia Graeca}, whose attribution of that work is the printing's.
The other witnesses were read in text layers, transcriptions or web texts; St John
Chrysostom's Greek exposition of Ps~113 is described and not quoted. Rupert of Deutz, Honorius Augustodunensis, Sicard of Cremona and William
Durandus are cited for what they say of the Introit, Collect, Epistle, Gradual,
Offertory and Communion, which the Mass they expounded shared with this formulary
although it read another Gospel; Rupert, Honorius and Durandus are cited on this
formulary's Gospel at the chapters where their books expound it. Bl.\ Ildefonso
Schuster and the continuation of \work{The Liturgical Year} comment on this
formulary as the 1962 Missal gives it. Origen is cited only for the view St Jerome
rejects and St Thomas names. The sentence St Rabanus gives under St Gregory's name
was not checked in Gregory's own works. Among the
witnesses not reached are St Ambrose on Ps~118, St Hilary and St Robert Bellarmine
on Ps~118:1, Greek commentators on Pss~89 and 118 other than the exposition named
above, Theophylact and Cornelius a Lapide on Ephesians, Hugh of Saint-Cher, St
Augustine's \work{De spiritu et littera}, Amalarius on the Offertory's verses, and St
Anthony of Padua's Sunday sermons.

\textbf{Chronology.} The rows of Scriptural Date and Location come from the
Scripture chronology corpus, through the record generated for this formulary: its
default profile's answer for every element and, for the Gospel and the Epistle, a
declared comparison under its critical profile. The record holds no occasion for any of the psalms, no
attribution for Ps~89 or for Job, and no date for the telling of the parable; it is
audited in \texttt{research/scope.md}.

\textbf{Rights.} The Latin of the formulary is published on the basis of its
public-domain antecedent, the Pustet Missal of 1862 (17 U.S.C. 103(b)). The Bibles,
the 1861 English, Migne's texts and the other editions and translations quoted from
the References are in the public domain in the United States. Cassiodorus and
St Thomas on Job are quoted briefly from public-domain Latin in web transcriptions.
The \work{New American Bible} introductions are summarised, not reproduced.

\textbf{Records.} The collation is in \texttt{propers/verified.md}, the occurrence
and rubrics in \texttt{research/context.md}, the search, loci, disagreements and
negative results in \texttt{research/scope.md}, and the reasoning behind the three
readings in \texttt{research/interpretations.md}; all were reviewed independently
before this study was written.
\par}

\phantomsection
\section*{References}
% A starred heading sets no mark, so the Appendix that ends above would go on
% heading the pages of the bibliography. The heading stays starred, because
% check-content-preflight finds the References section by that exact form; the
% mark is set explicitly through the leaf-local \runninghead, which the web
% converter's macro whitelist carries.
\runninghead{References}
\addcontentsline{toc}{section}{References}
\label{sec:references}

{\footnotesize
\subsection*{Liturgical books and Bibles}
\begin{itemize}
\item \work{Missale Romanum ex decreto Sacrosancti Concilii Tridentini restitutum, Summorum Pontificum cura recognitum}, editio typica (Typis Polyglottis Vaticanis, 1962), CMAA facsimile, \url{https://media.churchmusicassociation.org/pdf/missale62.pdf}: \work{Rubricae generales} 16, 91, 108, 111, 127; \work{Rubricae generales Missalis} 427, 434, 437, 444--446, 453, 475, 483, 494; \work{Ordo Missae} no.~1130; pp.~413--415 (nos.~1699--1708); \work{Tabula orationum}, p.~427 (nos.~1806--1808); pp.~689--690 (nos.~3883, 3889, 3891).
\item \work{Missale Romanum} (Ratisbon: Pustet, 1862), pp.~354--355, \latin{Dominica XXI. post Pentecosten}.
\item \work{The Roman Missal, Translated into the English Language for the Use of the Laity} (Philadelphia: Eugene Cummiskey, 1861), pp.~450--452 (XXI Sunday after Pentecost: Introit, Collect, Offertory, Secret, Communion, Postcommunion).
\item The Holy Bible, Douay--Rheims version, revised by Bishop Richard Challoner (Project Gutenberg text): 3~Kgs 6:1; Esth 11:1; 13:8--18; Job 1:1--22; 2:3--10; Ps 1:5; 89:1--4; 113:1--4; 118:1, 81--88; Lam 4:21; Mt 17:23; 18:1, 21--35; 20:2; Lk 6:37--38; 10:19; Acts 1:7; Rom 13:14; 16:20; 1~Cor 10:13; Eph 3:1; 4:1; 6:10--20; 1~Jn 1:8; Apoc 6:10--11.
\item \work{Biblia Sacra Vulgatae editionis} (Clementine Vulgate), eBible.org text: Esth 13:9--17; Job 1:1, 12; 2:6--7; Ps 89:1--2; 113:1; 118:1, 81--86; Mt 18:23--35; Eph 6:10--17.
\end{itemize}

\subsection*{Fathers, Doctors and saints}
\begin{itemize}
\item St Augustine, \work{Enarrationes in Psalmos}, trans. in Nicene and Post-Nicene Fathers, 1st ser., vol.~8 (CCEL text): on Ps~89:1--2; 113:1--4; 118, sermon~1 (vv.~1--3) and \latin{Caph} (vv.~81--88).
\item St Augustine, \work{Sermons} 58 and 83 (Benedictine numbering; Sermons 8 and 33 in the translation) and \work{De sermone Domini in monte} I.22.76--77, trans. in Nicene and Post-Nicene Fathers, 1st ser., vol.~6 (CCEL text).
\item St Gregory the Great, \work{Morals on the Book of Job} (\work{Moralia in Iob}), Library of the Fathers (Oxford): Preface 1--3, 7--8, 14, 16; II.19; IX.4--5.
\item St Jerome, \work{Commentarii in Matthaeum} III, on Mt 18:23--35, and \work{Commentarii in Epistolam ad Ephesios} III, on Eph 6:11--12, in Migne, PL 26 (Paris, 1845), cols.~132--133 and 543--546.
\item St Hilary of Poitiers, \work{Tractatus super Psalmos}, in Ps.~118, \latin{Caph}~5, and \work{Commentarius in Matthaeum} 18.10--11, in Migne, PL 9 (Paris, 1844), cols.~573--574 and 1023--1024.
\item St John Chrysostom, \work{Homilies on Matthew} 61, Nicene and Post-Nicene Fathers, 1st ser., vol.~10 (CCEL text); \work{Homilies on Ephesians} 22--24, Nicene and Post-Nicene Fathers, 1st ser., vol.~13 (CCEL text); \work{Expositio in Psalmos}, in Ps.~113, Migne, PG 55 (text layer; described).
\item St Thomas Aquinas, \work{Summa theologiae} I q.~103 a.~7, trans. Fathers of the English Dominican Province; \work{Super Evangelium S.~Matthaei lectura} c.~18, in \work{Opera}, editio altera Veneta, vol.~3 (Venice, 1745), pp.~242--244; \work{In Epistolam ad Ephesios} c.~6, lect.~3--4, in \work{In omnes D.~Pauli Apostoli epistolas commentaria}, vol.~2 (Li\`ege: H.~Dessain, 1857), pp.~356--359; \work{Expositio super Iob ad litteram} c.~1, Corpus Thomisticum web text.
\item St Rabanus Maurus, \work{Expositio in librum Esther}, preface, VI, XI and XIII, in Migne, PL 109 (Latin Wikisource transcription), cols.~635--636, 652, 664 and 668--669; \work{Commentaria in Matthaeum}, on Mt 18:23--35, in Migne, PL 107 (text layer).
\item St Robert Bellarmine, \work{A Commentary on the Book of Psalms}, trans. J.~O'Sullivan, eCatholic2000 transcription: on Pss~89:1--2; 113:1; 118:81--86; its prefaces I and II paraphrased, not quoted (Preface II, extant in manuscript only, trans. Michael J.~Miller).
\item St Alphonsus Liguori, \work{Sermons for All the Sundays in the Year}, trans. N.~Callan (Dublin: James Duffy, 1882), Sermon L, pp.~380--387.
\item Cassiodorus, \work{Expositio Psalmorum}, on Pss~89:1--2; 113:1; 118:1 and 81--86, Corpus Corporum web text.
\item \work{Expositiones in Psalmos}, printed under the name of St Athanasius, in Migne, PG 27 (Paris): on Ps~89:1--2, cols.~395--396; on Ps~118:81--86, cols.~493--496.
\item Origen, \work{Commentary on Matthew} XIV.10--11, trans. in Ante-Nicene Fathers, vol.~9 (New Advent text).
\end{itemize}

\subsection*{Liturgical commentators}
\begin{itemize}
\item Bl.\ Ildefonso Schuster, \work{The Sacramentary (Liber Sacramentorum)}, vol.~3 (London: Burns Oates \& Washbourne, 1927), pp.~178--181.
\item The continuation of Prosper Gu\'eranger, \work{The Liturgical Year} (Dom Lucien Fromage), \work{The Time after Pentecost}, series vol.~11, trans. Laurence Shepherd (Dublin: James Duffy and Sons, 1883), pp.~454--470.
\item Rupert of Deutz, \work{De divinis officiis} XII.21--22 and \work{De omnipotentia Dei} 9, in Migne, PL 170 (text layer).
\item Honorius Augustodunensis, \work{Gemma animae} IV.90--93, in Migne, PL 172 (Paris, 1854), cols.~724--726.
\item Sicard of Cremona, \work{Mitrale} VIII.21, in Migne, PL 213 (text layer).
\item William Durandus, \work{Rationale divinorum officiorum} VI.138--139, vol.~2 (Lyon, 1612; text layer).
\end{itemize}

\subsection*{Chronology sources}
\begin{itemize}
\item \work{The Catholic Encyclopedia} (New York, 1907--1912): ``Biblical Chronology'' (vol.~3, 1908); ``David'' (vol.~4, 1908); ``Esther'' and ``Epistle to the Ephesians'' (vol.~5, 1909); ``Gospel of Saint Matthew'' (vol.~10, 1911); ``St.\ Paul'' (vol.~11, 1911); ``The New Testament'' (vol.~14, 1912).
\item G.~L. Haydock, \work{The Douay--Rheims Bible with Haydock's Commentary}, note on Ps~118:1.
\item United States Conference of Catholic Bishops, \work{New American Bible, Revised Edition}, introductions to Job, the Psalms, Matthew and Ephesians, on their dating; summarised, not reproduced.
\end{itemize}
\par}
